\begin{lemma}[Forced Values Are Unique]
\label{lem:forced-values-are-unique}
Let $(P,S,1)$ be a Peano system, let $(W,c,g)$ be iterator data for it, and let
$R^*$ be the minimal iterator relation. If $(n,w_0)\in R^*$ and
$(n,w_1)\in R^*$, then $w_0=w_1$.
\hyperref[prf:forced-values-are-unique]{Go to proof.}
\end{lemma}
\end{lemmabox}
\LeanFormalizes{lem:forced-values-are-unique}{lra-lean}{LRA.VolumeII.PeanoSystems.Recursion.Iterator}{ForcedValuesAreUnique}{pending}

\begin{remark*}[Standard quantified statement]
\[
\forall n\in P\;\forall w_0,w_1\in W\;
\bigl((n,w_0)\in R^*\land (n,w_1)\in R^*\Longrightarrow w_0=w_1\bigr).
\]
\end{remark*}

\begin{remark*}[Predicate reading]
\[
\operatorname{FunctionalRelation}(R^*,P,W).
\]
\end{remark*}

\begin{remark*}[Negated quantified statement]
\[
\exists n\in P\;\exists w_0,w_1\in W\;
\bigl((n,w_0)\in R^*\land (n,w_1)\in R^*\land w_0\neq w_1\bigr).
\]
\end{remark*}

\begin{remark*}[Failure modes]
\begin{description}
  \item[Exposition.]
  Forced value uniqueness fails exactly when one stage is assigned two
  distinct values by the minimal iterator relation.
  \item[Two forced values.]
  \[
  \neg\operatorname{FunctionalRelation}(R^*,P,W).
  \]
  \[
  \exists n\in P\;\exists w_0,w_1\in W\;
  \bigl((n,w_0)\in R^*\land (n,w_1)\in R^*\land w_0\neq w_1\bigr).
  \]
\end{description}
\end{remark*}

\begin{remark*}[Contrapositive quantified statement]
\[
\exists n\in P\;\exists w_0,w_1\in W\;
\bigl((n,w_0)\in R^*\land (n,w_1)\in R^*\land w_0\neq w_1\bigr)
\Longrightarrow
R^* \text{ is not functional at every stage}.
\]
\end{remark*}

\begin{remark*}[Interpretation]
No stage can be forced to carry two different values. This is the key
determinacy step in the minimal-relation construction used to turn recursive
specifications into functions \citep{FefermanNumberSystems1964}.
\end{remark*}

\begin{dependencies}
\begin{itemize}
  \item \hyperref[lem:iterator-relation-consistency]{The Minimal Iterator Relation Is an Iterator Relation}
\end{itemize}
\end{dependencies}

